\documentclass[11pt]{article}
\usepackage{icmlko}
\begin{document}

\title{겹치는 축은 낱개로 재면 없는 것처럼 보인다 --- 수집은 묶음으로 값이 매겨진다}
\author{월드모델 실험실}
\date{노트 489 · 2026년 8월 3일}
\maketitle

\begin{abstract}
노트 488 이 ``격차는 축의 개수가 아니라 어느 수집이 붙었냐다''로 다시
썼다. 그러면 \textbf{수집마다 얼마인지} 물을 수 있다. 무너진 중앙값
나누기를 버리고 \textbf{수집의 있음/없음}으로 직접 갈라 여덟 짝을 쟀다
(양쪽 $25$ 행 이상 · $|$날 격차$| > 0.05$).

\textbf{먼저 계기가 또 한 번 틀렸다.} 축을 하나씩 가리면 몫 중앙값이
$14\%$ 라 ``자료가 붙은 행이 잘 맞는 건 그 자료 때문이 아니다''로 읽힌다.
그런데 \texttt{wiki\_level} $\cdot$ \texttt{momentum} $\cdot$
\texttt{volatility} 는 \textbf{한 수집의 축 셋}이라 하나를 가리면 나머지
둘이 대신한다. \textbf{수집 단위로 묶어 가리니 값이 나타났다.}

\begin{center}
\begin{tabular}{llr}
\hline
수집 & 낱개로 잰 몫 & 묶어 잰 몫 \\
\hline
게임 트렌드 & $2\% \cdot 45\% \cdot 29\%$ & $\mathbf{91\%}$ \\
게임 위키 & $23\% \cdot 28\% \cdot 14\%$ & $\mathbf{84\%}$ \\
세계애니 위키 & $1\% \cdot 6\% \cdot 87\%$ & $\mathbf{71\%}$ \\
\hline
\end{tabular}
\end{center}

세계애니 위키는 낱개 셋 중 둘이 \textbf{$1\%$ 와 $6\%$} 다 --- 그것만
보면 버려도 될 자료다. 묶으면 $71\%$ 다.
\end{abstract}

\begin{figure}[t]
\centering
\includegraphics[width=\textwidth]{figs/f1.pdf}
\caption{\textbf{가} --- 낱개 대 묶음. \textbf{나} --- 여덟 짝 전부.}
\end{figure}

\section{그리고 두 갈래로 갈린다}

묶어서 재면 여덟 짝이 깨끗이 둘로 나뉜다.

\begin{center}
\begin{tabular}{llrr}
\hline
수집 & 갈래 & 날 격차 & 몫 \\
\hline
게임 트렌드 & 행동 시계열 & $+0.0770$ & $\mathbf{91\%}$ \\
게임 위키 & 행동 시계열 & $+0.1119$ & $\mathbf{84\%}$ \\
세계애니 위키 & 행동 시계열 & $+0.1280$ & $\mathbf{71\%}$ \\
\hline
모바일 \texttt{gen} & 속성 표지 & $+0.0541$ & $27\%$ \\
애니 \texttt{gen} & 속성 표지 & $+0.1958$ & $14\%$ \\
세계애니 \texttt{goods\_scale} & 속성 표지 & $+0.1757$ & $1\%$ \\
만화 \texttt{goods\_scale} & 속성 표지 & $+0.0982$ & $-7\%$ \\
애니 \texttt{venue\_prominence} & 속성 표지 & $-0.0982$ & $-50\%$ \\
\hline
\end{tabular}
\end{center}

\textbf{행동 시계열 $3/3$ 이 $71\%$ 위 · 속성 표지 $5/5$ 가 $27\%$
아래다.} 겹치는 짝이 없다.

읽으면 --- \textbf{검색·조회 기록이 붙은 행이 잘 맞는 것은 그 기록
덕이다. 속성 표지가 붙은 행이 잘 맞는 것은 그 행이 원래 다른 종류라서다.}

\textbf{함정을 적는다.} 시계열 수집은 축이 셋이고 속성 수집은 하나다 ---
크기가 섞여 있다. 실제로 \emph{축 하나당}으로 보면 위키 낱개
($14 \sim 28\%$)와 속성 낱개($1 \sim 27\%$)가 같은 자리에 있다. 그러니
``시계열 축이 더 좋다''가 아니라 \textbf{``시계열은 묶음으로만 값이
있고 속성은 묶어도 없다''}가 맞는 말이다.

\section{애니 \texttt{venue\_prominence} 의 $-50\%$}

부호가 둘 다 뒤집힌 유일한 짝이다 --- 그 축이 \emph{있는} 행이 $0.0982$
\textbf{못 맞히는데}, 그 축을 가리면 $0.0495$ 더 못 맞힌다. 즉
\textbf{축은 돕는데 행이 어렵다.}

노트 486 이 애니 하나에서 찾은 것($89\%$ 가 라벨 압축)의 일반형이다.
그리고 노트 488 이 보였듯 그 부분 무리가 곧 애니의 ``축 많은 절반''
자체였다. \textbf{세 노트가 같은 $119$ 행을 세 번 말했다.}

\section*{조항 둘}

\begin{quote}
\textbf{겹치는 축을 낱개로 절제하면 그 수집이 쓸모없어 보인다.} 절제는
\textbf{수집 단위}로 한다. 세계애니 위키는 낱개로 $1\%$ 인데 묶으면
$71\%$ 다.
\end{quote}

\begin{quote}
\textbf{``자료가 붙은 행이 잘 맞는다''는 두 가지 뜻이다} --- 자료가 도운
것이거나(시계열), 그 자료가 붙는 행이 원래 다른 종류인 것이거나(속성).
\textbf{둘을 가르려면 그 자료를 가려 봐야 한다.} 덮음 표만 봐서는 못
가른다.
\end{quote}

\section*{남은 것}

\begin{itemize}
\item \textbf{수집 붙이기의 값어치} --- 시계열이 안 붙은 행에 붙이면
얼마 버나. 게임에서 위키 없는 $131$ 행이 $+0.5659$, 있는 $49$ 행이
$+0.6778$ 이다. 위 어림은 $+0.11$ 이지만 \textbf{붙는 행이 원래 다르니}
그대로 못 믿는다.
\item \textbf{속성 표지는 왜 묶어도 없나} --- \texttt{gen} 은 판에서
채택된 축인데 여기서는 $14 \sim 27\%$ 다. 판이 사는 것과 이 격차가
사는 것이 다르다.
\item \textbf{게임에서만 중앙값 나누기가 성립한다}(노트 488) --- 노트
$480 \sim 487$ 을 게임에서만 다시 묻는 일은 아직 안 했다.
\end{itemize}

\end{document}
